% Minimal filled monthly report used by the tests: same section structure as the local template
% (outputs/report/template-latex/template.tex), with generic text.
\documentclass[11pt,a4paper]{article}

\newcommand{\reporttitle}{Índice de Desconforto de Crédito}
\newcommand{\mesreferencia}{julho de 2026}
\newcommand{\mesref}{jul-2026}
\newcommand{\mesanterior}{jun-2026}
\newcommand{\competencia}{202609}
\newcommand{\mespublicacao}{setembro de 2026}
\newcommand{\proxdivulgacao}{outubro de 2026}
\newcommand{\mesproximo}{ago-2026}
\newcommand{\janelainicio}{jan-2014}
\newcommand{\seriesinicio}{mar-2011}
\newcommand{\reportsubtitle}{Nota Técnica de Atualização --- Divulgação \mespublicacao; competência \mesreferencia}
\newcommand{\reportdate}{Setembro de 2026}
\newcommand{\authorone}{Lauro Gonzalez}
\newcommand{\authortwo}{Matheus L. Carrijo}
\newcommand{\authorthree}{Rafael Schiozer}

\begin{document}

\section{Resumo}

Com a publicação, em \mespublicacao, das Estatísticas Monetárias e de Crédito do
Banco Central do Brasil (competência \competencia), o Índice de Desconforto
de Crédito (IDC) foi recalculado e estendido até \mesreferencia.

Em \textbf{\mesref}, o IDC foi de \textbf{0,963}, com alta de \textbf{0,013~ponto}
em relação aos \textbf{0,950} de \textbf{\mesanterior}.

\section{Resultados de \mesreferencia}

\begin{table}[htbp]
  \centering
  \caption{Componentes brutos e valores normalizados do IDC --- \mesref.}
  \label{tab:componentes}
  \begin{tabular}{p{8.0cm}rr}
    \toprule
    Indicador & Valor bruto & Valor normalizado \\
    \midrule
    IDC (média dos componentes normalizados)
      & ---
      & \textbf{0,963} \\
    C --- Comprometimento de renda das famílias com a dívida
      & 28,7\%
      & 1,000 \\
    I --- Inadimplência da carteira de crédito livre PF (90+ dias)
      & 7,8\%
      & 1,000 \\
    Q --- Participação das modalidades onerosas no crédito livre PF
      & 24,4\%
      & 0,888 \\
    \bottomrule
  \end{tabular}
\end{table}

\fonte{Banco Central do Brasil (Estatísticas Monetárias e de Crédito,
divulgação \competencia), elaboração própria.}

Em relação à observação imediatamente anterior --- \textbf{\mesanterior}, quando o IDC
marcava \textbf{0,950} --- o índice subiu \textbf{0,013~ponto} em \textbf{\mesref}.
O movimento concentrou-se em I.

\begin{itemize}
  \item Comprometimento de renda (C) permaneceu em \textbf{28,7\%} (\textbf{0,0~p.p.}).

  \item Inadimplência (I) avançou de \textbf{7,4\%} para \textbf{7,8\%} (\textbf{+0,4~p.p.}).

  \item Participação do crédito oneroso (Q) caiu de \textbf{24,5\%} para \textbf{24,4\%} (\textbf{-0,1~p.p.}).
\end{itemize}

Em conjunto, os dados indicam aumento do desconforto agregado em \textbf{\mesref}.

% SEÇÃO OPCIONAL — REVISÕES DOS DADOS
% Não inclua essa seção no relatório mensal por padrão.

\section{Trajetória do índice}

A Figura~\ref{fig:indice} mostra a evolução do IDC; a Figura~\ref{fig:componentes}
mostra as séries brutas de C, I e Q.
\section{Próxima atualização}

A próxima atualização é esperada para \proxdivulgacao; o IDC será estendido para
\mesproximo{} caso a série de comprometimento de renda (SGS~29034) avance até lá.

\section*{Notas}

\begingroup\footnotesize

Sobre o IDC. O Índice de Desconforto de Crédito (IDC) é construído a partir de
três séries do Sistema Gerenciador de Séries Temporais (SGS) do Banco Central do Brasil.

Revisões dos dados. Os valores históricos do IDC são recalculados a cada edição
com a safra mais recente dos dados do Banco Central. Revisões com efeito material
sobre o índice ou sua interpretação são descritas no relatório; a série
disponibilizada incorpora todas as revisões identificadas.

Citação recomendada. Gonzalez, L.; Carrijo, M.~L.; Schiozer, R.
\textit{Índice de Desconforto de Crédito (IDC).} FGVcemif/FGV-EAESP\@.
\endgroup

\clearpage
\section*{Anexo de figuras}

\begin{figure}[H]
  \centering
  \IfFileExists{index.png}{%
    \includegraphics[width=\linewidth]{index.png}%
  }{%
    \fbox{\rule{0pt}{14cm}%
      \rule{\dimexpr\linewidth-2\fboxsep-2\fboxrule\relax}{0pt}}%
  }
  \caption{Evolução do Índice de Desconforto de Crédito (\janelainicio{} a \mesref).}
  \label{fig:indice}
  \fonte{BCB, elaboração própria.}
\end{figure}

\clearpage
\begin{figure}[H]
  \centering
  \IfFileExists{components_raw.png}{%
    \includegraphics[width=\linewidth]{components_raw.png}%
  }{%
    \fbox{\rule{0pt}{17cm}%
      \rule{\dimexpr\linewidth-2\fboxsep-2\fboxrule\relax}{0pt}}%
  }
  \caption{Evolução das componentes brutas C, I e Q (\seriesinicio{} a \mesref).}
  \label{fig:componentes}
  \fonte{BCB, elaboração própria.}
\end{figure}

\end{document}
